\documentclass[12pt,a4paper]{article}

% --- Настройка языка и шрифтов для LuaLaTeX ---
\usepackage{fontspec}
\defaultfontfeatures{Ligatures=TeX,Scale=MatchLowercase}

\setmainfont{Alice-Regular}[
    Path=font/,
    Extension=.ttf,
    BoldFont = Alice-Regular,
    ItalicFont = Alice-Regular,
    BoldItalicFont = Alice-Regular,
    BoldFeatures = {FakeBold=1.6},
    ItalicFeatures = {FakeSlant=0.2},
    BoldItalicFeatures = {FakeBold=1.6,FakeSlant=0.2}
]

\newfontfamily\cyrillicfont{Alice-Regular}[
    Path=font/,
    Extension=.ttf,
    BoldFont = Alice-Regular,
    ItalicFont = Alice-Regular,
    BoldItalicFont = Alice-Regular,
    BoldFeatures = {FakeBold=1.6},
    ItalicFeatures = {FakeSlant=0.2},
    BoldItalicFeatures = {FakeBold=1.6,FakeSlant=0.2}
]

\usepackage{polyglossia}
\setmainlanguage{russian}
\setsansfont{TeX Gyre Heros}
\setmonofont{TeX Gyre Cursor}

\usepackage{graphicx}        					% Графика
\usepackage{amsmath, amsfonts, amssymb, amsthm}	% Математика
\usepackage{mathtools}
\usepackage{chemfig}
\usepackage[version=4,arrows=pgf-filled,
            textfontname=sffamily,
            mathfontname=mathsf]{mhchem}
\usepackage[protrusion=false]{microtype}       % Микротипографика
\usepackage{braket} 						% Braket нотация для квантовой механики
\usepackage{unicode-math} 					% Современный пакет для математики в LuaLaTeX
\usepackage{longtable}       				% Таблицы с разрывом страниц
\usepackage{tikz}           				% Рисование
\usepackage[europeanresistors]{circuitikz}	% Рисование электрических схем в европейском стиле
\usepackage{listings}						% Для отображения кода
\usepackage{tabularx}       				% Для таблиц во всю ширину
\usepackage{array}		  					% Для расширенных возможностей таблиц
\usepackage{cancel}                % Перечёркивать
\usepackage{float}          				% Для точного позиционирования фигур
\usepackage{pgfplots}                          % Для построения графиков
\usepackage{caption}                           % Для настройки подписей к рисункам и таблицам
\usepackage{xfp}                               % Для вычислений в LaTeX
\usepackage{luacode}                           % Для выполнения Lua-кода в документе    
\pgfplotsset{compat=1.18}
\usepackage[paper=letterpaper,
            top=20mm,
            bottom=20mm,
            left=10mm,
            right=10mm,
            footskip=0.4in 
]{geometry}
\usepackage{colortbl}

\usetikzlibrary{arrows.meta}

% --- Колонтитулы ---
\usepackage{fancyhdr}
\pagestyle{fancy}
\fancyhf{}

% --- Теоремы и окружения ---
\theoremstyle{plain}
\newtheorem{theorem}{Theorem}[section]
\newtheorem{lemma}[theorem]{Lemma}
\newtheorem{proposition}[theorem]{Proposition}
\newtheorem{corollary}[theorem]{Corollary}

\theoremstyle{definition}
\newtheorem{definition}[theorem]{Definition}

\theoremstyle{remark}
\newtheorem{remark}[theorem]{Remark}

\numberwithin{equation}{section}


\renewcommand{\headrulewidth}{0.4pt}
\renewcommand{\footrulewidth}{0.4pt}
\renewcommand\lstlistingname{Исходный код}

\lstdefinestyle{main}{
	backgroundcolor=\color{gray!8},   
	commentstyle=\color{green!50!black},
	keywordstyle=\color{blue!80!black},
	numberstyle=\small\color{darkgray},
	stringstyle=\rmfamily\color{orange!90!black},
	basicstyle=\ttfamily\small,
	breakatwhitespace=false,         
	breaklines=true,                   
	keepspaces=true,                 
	numbers=left,                    
	numbersep=8pt,                  
	showspaces=false,                
	showstringspaces=false,
	showtabs=false,                  
	tabsize=4
}	

\lstset{style=main}

\newcolumntype{C}{>{\centering\arraybackslash}X}
\renewcommand{\arraystretch}{1.25}

\begin{luacode*}
include = require("include")
\end{luacode*}

\begin{document}

\begin{center}
    \Large Лабораторная работа \number 2
\end{center}

\begin{center}
    \large Плазмотехническое травление в технологии интегральных схем
\end{center}

\textbf{Цель работы}: освоение технологии плазмохимического травления (ПХТ) материалов, применяемых в 
микроэлектронике; изучение факторов, влияющих на скорость и качество ПХТ.

\section{Теоретические сведения}

\subsection{Механизм ПХТ}

В основе процесса ПХТ лежит гетерогенная химическая реакция, проходящая на границе двух фаз: 
твердой и газообразной (или плазмообразной). Реакция протекает между поверхностными атомами 
обрабатываемого материала и химически активными частицами (ХАЧ).

\subsubsection{Химически активные частицы}

К химически активным частицам относятся радикалы — незаряженные частицы, имеющие на внешней 
оболочке неспаренные электроны и вледствие этого обладающие высокой химической активностью. 
Возникновение ХАЧ происходит в низкотемпературной плазме газового разряда, который может 
генерироваться как постоянным электрическим полем, так и переменными ВЧ- и СВЧ-полями.

В результате химической реакции между радикалами и поверхностными атомами 
обрабатываемого материала образуются продукты реакции, которые 
десорбируются с поверхности материала, переходят в газовую фазу и 
откачиваются вакуумным насосом.

\subsubsection{Стадии ПХТ}

Процесс ПХТ многостадийный: 
\begin{enumerate}
  \item доставка молекул газа в зону разряда;
  \item превращение молекул газа в радикалы под воздействием электронов разряда;
  \item доставка радикалов к поверхности материала, подвергаемого травлению;
  \item абсорбция радикалов на поверхности материала;
  \item химическая реакция радикалов с поверхностными атомами;
  \item десорбция продуктов реакции с поверхности материала;
  \item отвод продуктов реакции из реактора.
\end{enumerate}

\subsection{Скорость ПХТ}

Скорость ПХТ определяется скоростью наиболее медленной (лимитирующей) стадии.

\subsubsection{Стадия доставки молекул газа в зону разряда}

Определяется падением давления между газовым баллоном и 
зоной разряда реактора.

\subsubsection{Стадия диссоциации молекул газа под воздействием электронов разряда}

Скорость диссоциации молекул газа в электрическом разряде определяется 
числом столкновений молекул с электронами, энергия которых превышает 
некоторую пороговую энергию $E_\text{пор}$, необходимую для диссоциации молекул.
Притом нужно учитывать, что диссоциация молекул газа с необходимой пороговой
энергией $E_\text{пор}$ происходит с некоторой вероятностью 
\begin{equation*}
  P_\text{дис} = C_\text{г}(E_{e^-} - E_\text{пор}),
\end{equation*}
где $C_\text{г}$ — коэффициент пропорциональности (зависит от рода газа).

И тогда скорость диссоциации молекул газа (генерации ХАЧ)
\begin{equation*}
  V_\text{дис} = \left(\frac{8}{\pi m_{e^-}}\right)^{1/2} n_\text{г}n_{e^-} \frac{C_\text{г}(kT_{e^-})^{3/2}}{\overline{\lambda_{e^-}}} \left(2 + \frac{E_\text{пор}}{kT_{e^-}}\right)\exp{\left(-\frac{E_\text{пор}}{kT_{e^-}}\right)},
\end{equation*}
где $n_\text{г}$ — концентрация молекул газа, $n_{e^-}$ — концентрация электронов, $\overline{\lambda_{e^-}}$ — средняя длина свободного пробега электронов.

\subsubsection{Стадия доставки радикалов к поверхности материала, подвергаемого травлению}

Определяется температурой газа в ректоре, давлением газа и расстоянием от 
зоны разряда до поверхности материала.

\subsubsection{Стадия абсорбции радикалов на поверхности материала и 
стадия химической реакции радикалов с поверхностными атомами}

Скорость ПХТ возрастает с увеличением концентрации активных центров 
на поверхности кремния. Концентрация активных центров в свою очередь 
зависит от кристаллической структуры поверхности материала, а также 
от внешних воздействий (температуры, концентрации легирующих атомов примеси).

\subsection{Выбор газа для ПХТ}

При выборе газа учитывают скорость травления, селективность травления и 
анизотропность процесса.

Необходимым условием для травления материала ХАЧ является возможность 
образования летучих и стабильных при температуре процесса $T_\text{пр}$ 
(температура продуктов реакции). Летучесть можно оценивать по температуре
кипения-испарения продуктов реакции $T_\text{исп}$. Конкретно отношением 
$T_\text{исп}$ и $T_\text{пр}$:
\begin{enumerate}
  \item $T_\text{пр} << T_\text{исп}$ — травление невозможно (соединение нелетучее);
  \item $T_\text{пр} \approx T_\text{исп}$ — травление возможно, 
  скорость ограничена химической реакции, растёт с увеличением $T_\text{исп}$;
  \item $T_\text{пр} >> T_\text{исп}$ — травление возможно, 
  скорость ограничена диффузией реагента, скорость перестаёт расти.
\end{enumerate}

Необходимым условием является селективность травления 
\begin{equation*}
  S = \frac{V_\text{травл. материала}}{V_\text{травл. маскирующего материала}} \geq 1,
\end{equation*}

Необходимым условием является показатель анизотропии травления 
\begin{equation*}
  A = \frac{V_\text{травл. вглубь}}{V_\text{травл. в бок}} = \frac{L_\text{травл. вглубь}}{L_\text{травл. в бок}} \geq 1,
\end{equation*}

\subsection{Классификация ПХТ}

ПХТ бывает двух видов:
\begin{enumerate}
  \item Плазменное травление (ПТ) — травление происходит в области непосредственно плазмы газового разряда;
  \item Радикальное травление (РТ) — травление происходит в области, отделённой от плазмы газового разряда.
\end{enumerate}

Добавки в газ, которые при образовании соединений с молекулами газа образуют ХАЧ, 
ускоряют процесс ПХТ. Добавки в газ, которые образуют соединения с ХАЧ, замедляют
процесс ПХТ.

\subsection{Характеристика ПТ}

\subsubsection{Скорость ПТ}

Скорость ПТ описывается уравнением
\begin{equation*}
  V_\text{ПТ} = \frac{V_\text{хач}K_\text{тр}\tau_\text{ХАЧ}}{1 + \frac{bS_\text{мат}K_\text{тр}\tau_\text{ХАЧ}}{V_\text{плазмы}}},
\end{equation*}
где $V_\text{хач}$ — скорость генерации ХАЧ, 
$K_\text{тр}$ — константа скорости травления, 
$\tau_\text{ХАЧ}$ — время жизни ХАЧ, 
$b$ — константа, зависящая от материала подложки, 
$S_\text{мат}$ — площадь травления материала, 
$V_\text{плазмы}$ — объём зоны плазмы.

Если $S_\text{мат} << \frac{V_\text{плазмы}}{bK_\text{тр}\tau_\text{ХАЧ}}$, то скорость ПТ 
ограничена гетерогенной химической реакцией на поверхности материала и не зависит от 
площади материала.

Если $S_\text{мат} >> \frac{V_\text{плазмы}}{bK_\text{тр}\tau_\text{ХАЧ}}$, то скорость ПТ
ограничена диффузией ХАЧ к поверхности материала и обратно-пропорциональна площади материала.

Скорость ПТ существенно зависит от расхода газа с явным максимумом, от плотности мощности
разряда и от давления с явным максимумом.

\subsubsection{Анизотропия ПТ}

Описывается уравнением
\begin{equation*}
  L_\text{ширина линии в рабочей плёнке} = L_\text{ширина линии в маске} + \frac{2L_\text{толщина плёнки}}{A}\text{, где } A = \frac{L_\text{толщина плёнки}}{L_\text{бокового подтравливания}}.
\end{equation*}

\subsubsection{Селективность ПТ}

На селективность ПТ ключевым существенным образом влияют химические добавки в газ 
травления, при необходимой добавке селективность может быть повышена.

На селективность ПТ также влияет давление, в зависимости от условия протекания реакции.

Рост ВЧ-мощности понижает селективность ПТ, так как увеличивает скорость травления всех 
материалов.

Температура позволяет оставить один слой высоко-летучим, а другой низко-летучим, что
повышает селективность ПТ травления.

\subsection{Равномерность ПХТ}

Основным преимуществом РТ перед ПТ заключается в низком уровне радиационного и 
теплового воздействия на материал. В то же время РТ обладает высокой равномерностью
травления, что особенно важно при больших площадях подложек.

Для обесапечения высокой равномерности ПТ необходимо обеспечить однородность потока
плазмы по площади подложки.

\end{document}